\section{DO SACCHERIEGO Trzy typy płaszczyzn} % lang-pl
\label{sekcja_trzy_typy}

Saccheri dowodził swoich trzech hipotez ciągłością, a my wiemy dziś, że ciągłość jest zbędna. % lang-pl
W tej sekcji zobaczymy, jak jego wyniki wyglądają w dowolnej płaszczyźnie Hilberta, jak nazwiemy trzy przypadki i jak defekt trójkąta zacznie mierzyć odstępstwo od geometrii euklidesowej. % lang-pl

Hartshorne \cite[s. 307]{hartshorne_2000} zauważa, że Saccheri odróżni hipotezę kąta ostrego, prostego i rozwartego według tego, czy kąty przy górnej podstawie są ostre, proste czy rozwarte; udowodni, że jeśli któraś zachodzi dla jednego czworokąta, to dla wszystkich, ale użyje ciągłości (w postaci twierdzenia o wartości pośredniej). % lang-pl

\begin{theorem}
[Saccheriego] % lang-pl
\index{twierdzenie!Saccheriego (o trzech hipotezach)} % lang-pl
	W płaszczyźnie Hilberta, jeśli jeden czworokąt Saccheriego ma kąty przy górnej podstawie ostre, to wszystkie mają ostre; to samo dla kątów prostych i rozwartych. % lang-pl
\end{theorem}

Dowód (bez ciągłości, przez porównywanie czworokątów o równych środkowych): Hartshorne \cite[s. 307--309]{hartshorne_2000} (tw. 34.5, poprzedzone prop. 34.2--34.4). % lang-pl

\begin{theorem}
	W dowolnej płaszczyźnie Hilberta zachodzi jedna z trzech możliwości. % lang-pl
	Jeśli istnieje trójkąt o sumie kątów mniejszej od dwóch kątów prostych, to każdy trójkąt ma sumę mniejszą. % lang-pl
	Następujące warunki są równoważne: istnieje trójkąt o sumie kątów równej dwóm kątom prostym; istnieje prostokąt; każdy trójkąt ma sumę kątów równą dwóm kątom prostym. % lang-pl
	Jeśli istnieje trójkąt o sumie większej od dwóch kątów prostych, to każdy trójkąt ma sumę większą. % lang-pl
\end{theorem}

Dowód: Hartshorne \cite[s. 309--311]{hartshorne_2000} (tw. 34.7), na podstawie prop. 34.6: dla każdego trójkąta istnieje czworokąt Saccheriego, którego suma dwóch kątów przy górnej podstawie jest równa sumie kątów trójkąta \cite[s. 310]{hartshorne_2000}. % lang-pl

Te trzy przypadki Hartshorne nazywa geometrią \emph{półhiperboliczną} (suma mniejsza od $\pi$), \emph{półeuklidesową} (równa $\pi$) i \emph{półeliptyczną} (większa od $\pi$); słowo „hiperboliczna'' rezerwuje dla geometrii spełniającej aksjomat Hilberta o równoległych granicznych (zobacz sekcję \ref{sekcja_hiperboliczna}), bo półhiperboliczne płaszczyzny bez tego aksjomatu też istnieją \cite[s. 311]{hartshorne_2000}. % lang-pl
Przypadek półeliptyczny, w którym Saccheri odrzucił hipotezę kąta rozwartego, odkryje w 1900 roku Dehn, nazywając takie geometrie nielegendrowskimi; przykład to mały kawałek sfery nad ciałem nieArchimedesowym \cite[s. 311, 318]{hartshorne_2000}. % lang-pl
\index[persons]{Dehn, Max}%
\index{geometria!półhiperboliczna} % lang-pl

\begin{definition}
[defekt trójkąta] % lang-pl
	Defektem trójkąta nazywamy liczbę $\delta = 2\mathrm{RA} - (\alpha + \beta + \gamma)$, gdzie RA oznacza kąt prosty; jest ona zerem dla trójkąta euklidesowego, dodatnia w płaszczyźnie półhiperbolicznej i ujemna w półeliptycznej. % lang-pl
\end{definition}

Hartshorne \cite[s. 311]{hartshorne_2000}. % lang-pl

\begin{proposition}
	Jeśli odcinek z wierzchołka dzieli trójkąt na dwa trójkąty, to defekt całości jest sumą defektów części: $\delta(ABC) = \delta(ABD) + \delta(BCD)$. % lang-pl
\end{proposition}

Dowód: Hartshorne \cite[s. 311--312]{hartshorne_2000} (lemat 34.8). % lang-pl

Z tego wprost wyniknie pole trójkąta jako defekt (zobacz sekcję \ref{sekcja_suma_katow}), a w sekcji o polu nieeuklidesowym Hartshorne \cite[s. 326 nn.]{hartshorne_2000} rozwija to systematycznie. % lang-pl

Zadania: czworokąt Lamberta w trzech geometriach (zad. 34.2), kąt wpisany w półokrąg (zad. 34.3) i przystawanie trójkątów o równych kątach (zad. 34.4) -- Hartshorne \cite[s. 316]{hartshorne_2000}. % lang-pl
% Źródła: Wikipedia: Ferdinand Karl Schweikart, Franz Taurinus, Non-Euclidean geometry.
% https://en.wikipedia.org/wiki/Ferdinand_Karl_Schweikart
% https://en.wikipedia.org/wiki/Franz_Taurinus

